\let\negmedspace\undefined
\let\negthickspace\undefined
\documentclass[journal,12pt,onecolumn]{IEEEtran}
\usepackage{cite}
\usepackage{amsmath,amssymb,amsfonts,amsthm}
\usepackage{algorithmic}
\usepackage{graphicx}
\graphicspath{{./figs/}}
\usepackage{textcomp}
\usepackage{xcolor}
\usepackage{txfonts}
\usepackage{listings}
\usepackage{enumitem}
\usepackage{mathtools}
\usepackage{gensymb}
\usepackage{comment}
\usepackage{caption}
\usepackage[breaklinks=true]{hyperref}
\usepackage{tkz-euclide} 
\usepackage{listings}
\usepackage{gvv}                                        
\usepackage{amsmath}                                 
\usepackage[latin1]{inputenc}     
\usepackage{xparse}
\usepackage{color}                                            
\usepackage{array}                                            
\usepackage{longtable}                                       
\usepackage{calc}                                             
\usepackage{multirow}
\usepackage{multicol}
\usepackage{hhline}                                           
\usepackage{ifthen}                                           
\usepackage{lscape}
\usepackage{tabularx}
\usepackage{array}
\usepackage{float}

\begin{document}

\title{2.10.27}
\author{AI25BTECH11038 -- Tejas Uppala}
{\let\newpage\relax\maketitle}

\textbf{Question:}

\noindent The scalar $\textbf{A}\cdot\brak{\brak{\textbf{B} + \textbf{C}}\times\brak{\textbf{A} + \textbf{B} + \textbf{C}}}$ equals

\begin{multicols}{2}
\begin{enumerate}
    \item $0$
    \item $[\textbf{A} \textbf{B} \textbf{C}]$
    \item $[\textbf{A} \textbf{B} \textbf{C}] + [\textbf{B} \textbf{C} \textbf{A}]$
    \item None of these
\end{enumerate}
\end{multicols}

\textbf{Solution:}

\noindent The expression is a scalar triple product, which can be written using the box product notation:
\begin{equation}
E = \mathbf{A} \cdot ((\mathbf{B} + \mathbf{C}) \times (\mathbf{A} + \mathbf{B} + \mathbf{C})) = [\mathbf{A} \quad (\mathbf{B} + \mathbf{C}) \quad (\mathbf{A} + \mathbf{B} + \mathbf{C})]    
\end{equation}

\noindent The box product is equivalent to the determinant of the matrix formed by the vectors. We can use determinant properties (where rows are the vectors $\mathbf{R}_1, \mathbf{R}_2, \mathbf{R}_3$) to simplify:

\textbf{Row Operation:} Apply the determinant property that allows replacing one row by the sum/difference with other rows.

\textbf{Operation:} Replace the third vector ($\mathbf{R}_3$) by subtracting the first two vectors ($\mathbf{R}_1$ and $\mathbf{R}_2$):

\begin{equation}
    \mathbf{R}_3' = \mathbf{R}_3 - \mathbf{R}_1 - \mathbf{R}_2
\end{equation}

The new third vector is:
\begin{equation}
 \mathbf{R}_3' = (\mathbf{A} + \mathbf{B} + \mathbf{C}) - \mathbf{A} - (\mathbf{B} + \mathbf{C}) = \mathbf{A} + \mathbf{B} + \mathbf{C} - \mathbf{A} - \mathbf{B} - \mathbf{C}  \mathbf{0}   
\end{equation}
The box product becomes:
\begin{equation}
    E = [\mathbf{A} \quad (\mathbf{B} + \mathbf{C}) \quad \mathbf{0}]
\end{equation}

The determinant (or scalar triple product) is \textbf{zero} if one of the vectors is the zero vector ($\mathbf{0}$).
$$E = 0$$
The answer is \textbf{(A) 0}

\end{document}